\section{Ring-class switches, determinant lattices, and transfer}
\label{rd:section}

This section constructs the derivative classes and proves the determinant
and transfer statements.  The arithmetic models, the character \(\Psi\), and their
evaluation maps are those of Lemmas~\ref{rl:finite-model},
\ref{rl:finite-limits}, \ref{rl:evaluation}, and \ref{rl:kernel},
and Propositions~\ref{rl:characters} and \ref{rl:outer-evaluation}.
There are two outputs.  Corollary~\ref{rd:ring-lower} bounds a
Heegner-point index at a known simple zero, uniformly in the conductor
and with dependence only on a bound for the number of companion primes.
Theorem~\ref{rd:missing-vertex} instead fixes the companion field and
recovers an unknown base derivative from a family of known simple zeros.
The first step is an integral switch: adding a derivative prime replaces
a finite local condition by its singular counterpart without changing
the determinant valuation.  Repeated switches then give the uniform
lower bound.  For the missing vertex we need more: the determinants
must occur in one group-ring family with a uniformly bounded clearing
factor.  We construct the arithmetic diagrams before applying these two
arguments, and verify their specialization and family properties where
they are used.

Fix \(E/\mathbb Q\), put \(T=T_2(E)\), and normalize \(v(2)=1\).  Let
\[
 \Lambda=\mathbb Z_2[[t]],\qquad
 \mathcal O=\widehat{\Lambda_{(2)}},\qquad
 \Omega=\mathbb F_2((t)).
\]
For a perfect cohomological complex \(C\), use the inverse determinant line
\begin{equation}
 \mathcal D(C)=\bigotimes_i(\det C^i)^{(-1)^{i+1}}.
 \label{rd:det-convention}
\end{equation}
In particular, for a complex in degrees one and two this is
\(\det C^1\otimes(\det C^2)^{-1}\).  A basis of this integral line is fixed
whenever a determinant coordinate is mentioned.  Changing it multiplies
the coordinate by an integral unit.

The ring data consist of an imaginary quadratic field \(K=\mathbb Q(\sqrt k)\)
split at \(2N\), a product \(h\) of signed good odd primes split in \(K\),
and \((h,k)=1\).  The fields with extra units and the CM field of \(E\)
are excluded.  Write
\begin{equation}
 t_p=\dim_{\mathbb F_2} E[2]^{\operatorname{Fr}_p},\qquad
 w(h)=\frac12\sum_{p\mid h}t_p.
\label{rd:weights}
\end{equation}
Choose the modular parametrization \(\pi:X_0(N)\to E\)
normalized by \(\pi(\infty)=0\), translating it by the rational
point \(-\pi(\infty)\) first if necessary.  By the Manin--Drinfeld theorem
\cite{Manin1972,Drinfeld1973}, images of cusp
differences are torsion, so one fixed integer kills every
cuspidal ambiguity.  No assertion that an arbitrary translate
has torsion value at \(\infty\) is used.
The scalar deformation is trivial at \(2N\), has unit exponent on the
auxiliary conductor inertia, and odd exponent at the indicated split
conductor Frobenii.  Its auxiliary prime \(r\) is allowed to change with
the precision.  The trace bound in Proposition~\ref{rl:characters} is part of
the data.

\subsection{The local switch and its integral determinant}
\label{rd:local-subsection}

\begin{lemma}[The split local model]\label{rd:local-model}
Let a sequence of inert good primes \(\ell\) have rational Frobenius
lifts tending on \(T\) to a fixed complex conjugation \(J\).
At the unique prime of \(K\) above \(\ell\), the limiting local complex
has
\begin{equation}
 H^0=M,\qquad H^1=F\oplus S=M\oplus M,\qquad H^2=M,
 \label{rd:local-cohomology}
\end{equation}
with zero differentials in a free split model.  Here \(M=T\otimes\Lambda\)
with the relevant scalar coefficient structure.  The two degree-one
coordinates of a cocycle are
\begin{equation}
 f(c)=c(\gamma^2),\qquad s(c)=c(\sigma),
 \label{rd:coordinates}
\end{equation}
where \(\gamma\) is the chosen rational Frobenius lift and
\(\sigma\) is an oriented pro-\(2\) tame generator.  The plane
conditions \(F=\{s=0\}\) and \(S=\{f=0\}\) include \(H^0\) and exclude
\(H^2\).  Each is its own orthogonal complement for conjugate-Weil
duality.  The strict condition \(H^0\) and the relaxed condition
\(H^0\oplus H^1[-1]\) are mutually dual.  All four inclusions can be
equipped with compatible integral duality triangles.
\end{lemma}

\begin{proof}
The prime-to-\(\ell\) local cochains are computed by the tame inertia and
residue-field directions.  In these two directions the differentials
are formed from inertia minus one and Frobenius minus one, with the
usual Frobenius action on tame inertia.  On \(M\), inertia is trivial at
the derivative prime.  The square of a rational Frobenius is trivial
on every ring-class character.  In the limit it is also the identity on
\(T\), while \(\ell^2\) tends to one.  Thus both differentials become zero,
giving \eqref{rd:local-cohomology}.

Here is a finite-precision check of the pairing, needed at \(2\).
At a retained quotient killed by \(2^a\), impose also
\(q=\ell^2\equiv1\pmod{2^{a+1}}\), as well as trivial
Frobenius-square action on the coefficient module at this quotient.
The tame two-generator complex then has zero differentials,
with the coordinates in \eqref{rd:coordinates}.  Its degree-one
scalar generators are the unramified and tame characters.
More explicitly, use the presentation
\(\phi\sigma\phi^{-1}\sigma^{-q}=1\).
After augmentation with trivial coefficient action its Fox
resolution has terms of ranks \(1,2,1\), with
\(d^0=0\) and \(d^1=(0,1-q)\), in the ordered finite and
tame coordinates.  The diagonal from the same resolution
has zero pure-finite term, unit mixed terms, and pure-tame
coefficient \(\pm\binom q2\).  The latter coefficient is
the augmented second derivative of \(\sigma^{-q}\)
together with its adjacent \(\sigma\) in the relation.
Both \(1-q\) and \(\binom q2\) vanish modulo \(2^a\)
under the imposed congruence.  Thus this is an actual
split complex and an actual pure-plane zero pairing
at the retained precision.
The field contains the \(2^{a+1}\)-st roots of unity, so
\(-1\) is a \(2^a\)-th power.  The norm-residue identity
\((x,x)=(x,-1)\) for the \(2^a\)-Hilbert symbol therefore
makes the square of each scalar generator zero.
The two generators form a basis; local Tate duality makes
their cross pairing a unit.  The same calculation applies
coefficientwise to the free Artin coefficient module, including
its nilpotent \(t\) and quadratic variables.  Thus the pure
plane restrictions are already zero at every retained quotient.
All these congruences can be imposed simultaneously with the
derivative Frobenius conditions by increasing their precision.
Taking the compatible limit proves integral vanishing of the
pure squares.  It is not deduced from skew symmetry in
characteristic two.

Equivalently, in the limiting tame presentation
\(\phi\sigma\phi^{-1}=\sigma^q\) with \(q\to1\), the
rank-\((1,2,1)\) resolution specializes to the exterior
resolution for two procyclic directions.  The finite
Hilbert-symbol calculation specifies the cup product on
its split cohomology and rules out a hidden square term.
Conjugation transports the finite Tate coordinate by \(J\) and the
singular coordinate by \(-J\), with possible scalar units coming from
the orientations.  Consequently the cross pairing is
\begin{equation}
 \beta(x,y)=\epsilon\,\langle x,Jy\rangle,\qquad
 \epsilon\in\Lambda^\times,\quad x\in F,\ y\in S,
 \label{rd:cross-pairing}
\end{equation}
where \(\langle\, ,\,\rangle\) is the perfect alternating Weil form.
This proves the assertions about orthogonal complements.  Since the
restrictions of this explicit pairing to either pure plane are zero as
maps of the split complexes, their isotropy homotopies may be chosen to
be zero.  Only a product of two degree-one terms could map to the
degree-two local invariant target; the degree-zero terms introduce
no additional restriction.  These zero homotopies agree on the
strict subcomplex.  These are the actual
orthogonality data required in the Selmer duality construction
\cite[Theorem~6.3.4]{Nekovar2006}; the compatibility assertion is therefore
an assertion about these displayed complexes and maps.  Lemma~
\ref{rl:finite-limits} is used to retain this finite diagram together
with the unchanged local data.
\end{proof}

We now construct global classes with these local coordinates, using
the CM points \(y_m\in E(H_m)\) fixed in Section~\ref{rl:section}.
For the fixed twist put
\(c=|h|\).  If \(I\) is a finite set of derivative primes, put
\(n=\prod_{\ell\in I}\ell\); the points used for this set have
conductor \(crn\).

Let \(C_I\) denote the limiting Selmer complex with coefficients
\(T\chi_h\Psi\).  Its local conditions are Kummer at \(2N\), full
at the active primes dividing \(h\) and at \(r\), singular at the
primes in \(I\), and unramified elsewhere.  At a new derivative
prime \(\ell\notin I\), the unramified condition is the finite
plane \(F_\ell\) of Lemma~\ref{rd:local-model}.  Both plane
conditions retain local \(H^0\) and omit \(H^2\), as in that
lemma.  The next result constructs the class \(P_I\) in this
complex; its two displayed localizations are taken at the newly added
prime \(\ell\).

\begin{lemma}[Descent and the switch at \(2\)]\label{rd:derivative}
Perform each finite ring-class construction and its weighted transfer
before retaining its target-side data by the construction following
Lemma~\ref{rl:finite-limits}.  There is one
integer \(L\ne0\), depending on the fixed curve, parametrization and
fixed local types, such that the classes obtained from the Heegner
derivatives of \(L\) times the point satisfy, for every finite switched
set \(I\),
\begin{equation}
 P_I\in H^1(C_I),\qquad
 f(P_{I\cup\{\ell\}})=0,\qquad
 s(P_{I\cup\{\ell\}})=J f(P_I).
 \label{rd:switch}
\end{equation}
The same \(L\) works for every number of active and derivative primes.
The previously imposed singular conditions are retained.
\end{lemma}

\begin{proof}
We use Kolyvagin's derivative operators \cite{Kolyvagin1990}, keeping
their integral descent explicit at $2$.  Choose the orientations
coherently.  The relative ring-class group at
an inert prime \(\ell\) is cyclic of order \(\ell+1\).  For distinct new
primes these groups form a product: the ring-class exact sequence has
kernel the product of the local residue-unit quotients, and the only
global units here are \(1,-1\), whose images in these quotients are
trivial.  Choose a generator \(\sigma_\ell\) compatibly with the
tame orientation in Lemma~\ref{rd:local-model}; its pro-\(2\)
component gives the coordinate \(s\) at the retained precision.
Set
\[
 D_\ell=\sum_{i=1}^{\ell} i\sigma_\ell^i,\qquad
 D_I=\prod_{\ell\in I}D_\ell.
\]
Direct multiplication gives
\begin{equation}
 (\sigma_\ell-1)D_\ell=(\ell+1)-
             \sum_{i=0}^{\ell}\sigma_\ell^i.
 \label{rd:derivative-identity}
\end{equation}
The Hecke correspondence and reduction of its CM isogenies give, for
every class-group translate,
\begin{equation}
 \operatorname{Tr}(y_{c r n})=a_\ell y_{c r n/\ell},
 \qquad
 \widetilde y_{c r n}=\operatorname{Fr}_\ell
                       \widetilde y_{c r n/\ell}.
 \label{rd:heegner-relations}
\end{equation}
These are the point relations, not an application of a residual
irreducibility theorem.  A fixed multiple of the parametrization removes
cuspidal terms.

There is a uniform bound on the \(2\)-primary torsion over the union of
the ring-class fields in use.  To see this, choose a fixed nonidentity
scalar homothety in the \(2\)-adic Tate image.  Square a lift so that it
belongs to every quadratic subgroup under consideration.  Multiply it
by its conjugate under complex conjugation.  Its action on each
ring-class field is trivial, because conjugation inverts the abelian
ring-class action, while its action on \(T\) is still a fixed
nonidentity scalar.  Write that scalar minus one as \(2^e\) times
a unit.  Then \(2^e\) kills the torsion in question.  The existence of
this fixed scalar, including the CM case, is the Tate-image input
used in Lemma~\ref{rl:kernel}.

Work first with coefficients modulo \(2^a\).  Require \(a_\ell\) and
\(\ell+1\) to vanish to a precision substantially larger than \(a\).
Equations \eqref{rd:derivative-identity} and
\eqref{rd:heegner-relations} then make the derivative Kummer class
invariant under the relative ring-class group.  In the
inflation--restriction sequence, both the obstruction to descent and
the ambiguity of descent are killed by \(2^e\), since their
coefficient module is killed by \(2^e\).  Multiplication by \(2^{2e}\)
therefore gives a descent independent of this ambiguity: descend
\(2^e\) times the class and multiply its lift by \(2^e\).
There is no repetition of this factor for individual primes; descent
is performed once from the field with the complete conductor.

The target coefficient module is obtained by an integral transfer.
For a finite Artin coefficient ring \(A\), let
\(\Delta=\operatorname{Gal}(H_{cr}/K)\), through which the
finite scalar and quadratic characters factor.  Shapiro's
cochain map and the coefficient homomorphism
\[
 A[\Delta]\longrightarrow A,\qquad
 [\sigma]\longmapsto\chi_h(\sigma)\Psi(\sigma)
 \]
send the descended class to coefficients \(T\chi_h\Psi\).
For a universal quadratic family, replace the quadratic
value by its group element in \(A[G]\), also including the
fixed base character.  Choose inverse actions on both
sides when using that Shapiro convention.  Explicitly,
these maps are finite coset sums with the indicated scalar
weights; there is no division by \(\#\Delta\).
These chain maps, in all degrees, are used at this finite stage.
Naturality of restriction and Kummer corestriction supplies their local
lifts and comparison homotopies.  Apply the maps to the classes and
these witnesses before passing to bounded target models.  Neither
the complex over \(H_{cr}\) nor its induced coefficient module
\(A[\Delta]\) is retained in the limiting diagram.  We check the
required target local conditions next.

At a place over \(2N\) these descent extensions are unramified.
For the identity component of a Neron model, the obstruction to
unramified descent vanishes.  Indeed, Lang's theorem gives the assertion
on the connected special fiber, and the successive formal-group
quotients are additive residue-field modules, whose unramified cyclic
cohomology vanishes by the normal-basis and trace calculation.
Passing to the inverse limit of these quotients preserves the assertion.
The remaining obstruction is killed by the geometric component-group
order.  Those orders have a uniform bound over unramified extensions
of the finitely many fixed local twisting types.  Include their product,
the descent factor, and the parametrization factor in \(L\).
Thus the descended class has the integral Kummer condition at \(2N\).
No good-ordinary or semistability condition was used.  Transfer with
the scalar character preserves this condition because \(\Psi=1\)
there.

Here is the local calculation, including the division by two.  Apply
the derivatives at all primes other than \(\ell\), and write \(X,Y\)
for the resulting lower and upper points.  Put \(F_\ell=\operatorname{Fr}_\ell\).
The lower field splits at the \(K\)-prime \((\ell)\), since that ideal
is principal.  Reduction therefore gives \(F_\ell^2\widetilde X=\widetilde X\).
Good reduction identifies the prime-to-\(\ell\) torsion with the
torsion of the reduced curve.  In the following coordinate formulas,
\(F_\ell\) acts on reduced points, where it satisfies its
Frobenius polynomial.  If \(Q\) is a \(2^a\)-division point
of \(LX\), the finite coordinate
of its Kummer class is
\begin{equation}
 f=(F_\ell^2-1)\widetilde Q=
       \bigl(a_\ell F_\ell-(\ell+1)\bigr)\widetilde Q.
 \label{rd:finite-formula}
\end{equation}
We give the descent-cocycle calculation of the other coordinate.
Choose \(Z\) with \(2^aZ=LD_\ell Y\), and choose the descended
cocycle \(b\) so that on the upper-field Galois group it is
\(b(g)=(g-1)Z\).  A change by the already bounded descent ambiguity
has no effect on the eventual limit.  The crossed identity gives
\[
 U=(\sigma_\ell-1)Z-b(\sigma_\ell)
                  \in E(H_{crn}).
 \]
For \(g\) in the upper-field Galois group, compare the two
expansions of \(b(g\sigma_\ell)\); they give \(gU=U\).
Moreover,
\[
 2^aU=L\bigl((\ell+1)Y-a_\ell X\bigr).
 \]
Both integers \(\ell+1,a_\ell\) are divisible by \(2^a\).
Consequently
\[
 U_0=L\frac{\ell+1}{2^a}Y-
                      L\frac{a_\ell}{2^a}X
 \]
is an actual upper-field point with the same \(2^a\)-multiple
as \(U\).  Thus \(U-U_0\) is upper-field torsion of bounded
exponent.  This comparison with an actual integer quotient
is what bounds the ambiguity; arbitrary division points
would not give this bound.  Reduction makes
\(\sigma_\ell\) act trivially and
\(\widetilde Y=F_\ell\widetilde X\).  Therefore
reduction of \(b(\sigma_\ell)=-U\) gives its singular coordinate
\begin{equation}
 s=\bigl(a_\ell-(\ell+1)F_\ell\bigr)\widetilde Q
 \label{rd:singular-formula}
\end{equation}
up to the uniformly bounded descent torsion already described.
The Frobenius polynomial yields the exact operator identity
\begin{equation}
 (F_\ell-a_\ell)
     \bigl(a_\ell F_\ell-(\ell+1)\bigr)
       =a_\ell-(\ell+1)F_\ell.
 \label{rd:exact-frobenius}
\end{equation}
Thus \(s=(F_\ell-a_\ell)f\); when \(a_\ell\) vanishes at the
retained torsion precision, this is \(s=F_\ell f\).  Taking the
Frobenius limit gives \(s=Jf\).

For the new finite coordinate, reduction of the derivative point is
\begin{equation}
 \sum_{i=1}^{\ell}iF_\ell\widetilde X
       =\frac{\ell(\ell+1)}2F_\ell\widetilde X.
 \label{rd:half-factor}
\end{equation}
The conductor extension at \(\ell\) is totally ramified, so a
Frobenius lift may be chosen in the upper-field local Galois group.
There the descended cocycle restricts to the Kummer cocycle of the
derivative point.  If \(v_2(\ell+1)\ge a+1\), its displayed
reduction is \(2^a\) times an \(F_\ell^2\)-fixed point.
Its finite Kummer coordinate is therefore zero modulo \(2^a\).
The additional \(1\) is essential at
\(2\).  It is imposed in the choice of the next precision, and does
not multiply the class by a new denominator.  Bounded torsion
ambiguities disappear on passing from \(E[2^a]\) to \(E[2^m]\) with
\(a-m\ge e\).  Choose \(a\) and the Frobenius precision successively
faster than \(m\).

The calculation is compatible with every translate in the transfer:
transport the place and its Frobenius before applying that translate.
The scalar weight on the lower and upper class is identical.  Applying
the same calculation at each \(\ell\in I\) proves the earlier
singular conditions as well.

Finally complete the weighted transfer at each stage and retain the
resulting target Selmer-cone cycles and their local lifts.  An equality
of the transferred local classes gives a boundary witness in the
target local complex; retain that witness after contraction as well.
For instance, retain the degree-zero vectors whose boundaries are the
differences in the two switch identities.  Include the split local Fox
complexes, their zero isotropy homotopies, and the target comparison
matrices.  The target-side construction in Section~\ref{rl:section} proves
that this is a bounded diagram for each fixed active support and
finite switched set.  Its bound is independent of the degrees of
the ring-class fields.
When a further derivative prime is selected, extend the
same original finite-stage diagram and its contraction
maps, using the same inner ultrafilter.  Do not construct
an unrelated cohomologically equivalent old complex.
The earlier evaluations, including all degree-two
cocycle identities, are consequently retained by the
same contraction maps after each new switch.  Lemma~
\ref{rl:finite-limits} applies to these extended diagrams.
Their limiting classes give
\eqref{rd:switch}.  The displayed local formulas agree with the
calculation underlying \cite[Proposition~1.7.4]{Howard2012}; Howard's
odd-prime global hypotheses are not being asserted at \(2\).
\end{proof}

\paragraph{The paired arithmetic determinant.}
All the target classes and comparison maps are constructed in the same
finite diagram.
At each Artin precision, use the character of
Proposition~\ref{rl:characters}, the global Selmer cone, and the
Kummer complexes at \(2N\).  At a derivative prime impose enough
Frobenius precision for Lemma~\ref{rd:derivative} and also
\(\ell^2\equiv1\pmod{2^{a+1}}\) at every retained quotient killed
by \(2^a\).  Lemma~\ref{rd:local-model} then gives the finite,
singular, strict and relaxed inclusions with their common zero
isotropy homotopies.  Applying the actual global duality morphism
\eqref{rl:global-duality} constructs their integral localization
and duality triangles.  Perform the descent and weighted transfers of
Lemma~\ref{rd:derivative} at each finite stage.  Include their target
cycles, local Kummer lifts and target boundary witnesses in this
diagram before applying Lemmas~\ref{rl:finite-model} and
\ref{rl:finite-limits}, using only the target-side data.
All remaining maps have bounded target models as both source and
target; the ring-class source complexes are not among these objects.

A further switch extends these original finite stages, their
contractions and their fixed inner ultrafilter.  Theorem~
\ref{rl:limit-evaluation-package} therefore retains the same
degree-zero, degree-one and degree-two evaluations and the
exact-kernel restrictions.  In particular, the old evaluations
are not replaced by maps on an independently chosen equivalent
complex.

Here is how a class produces its paired functional.  Work over a
height-one localization where the full local complexes at the active
primes and \(r\) are acyclic.  By Lemma~\ref{rl:local-acyclic}
this includes \(\mathcal O\) and every characteristic-zero test
used below.  These complexes have integral contractions.  They
identify the relaxed Selmer complex with its orthogonal complex;
at switched primes the plane condition is already self-dual, and
at \(2N\) the integral Kummer condition is self-dual.  Lift the
class through this comparison and compose it with the retained
global duality map.  The result is a closed degree-two functional
\(\lambda_I\), integral up to the single initial multiplier and
any fixed old-support correction.  Thus it is not an arbitrarily
chosen rational cohomology dual.

When the generic cohomology has dimension one in degrees one and
two and is zero elsewhere, define \(u_I\) to be the coordinate of
\[
 [P_I]\otimes[\lambda_I]\in
 H^1(C_I)\otimes H^2(C_I)^*
       \simeq\mathcal D(C_I)
\]
over the fraction field, in the fixed integral determinant basis.
For an acyclic constituent set this class-dual coordinate to zero.
The generic functional is determined by the same derived duality
at all localizations; different local contractions represent that
functional, not one cochain asserted to be integral everywhere.
For \(I=\varnothing\) write \(u_h\).  The rank-one assertions
needed for the arithmetic applications are verified below by
rational central specialization.

The arithmetic switch must now be compared in these determinant
lattices.  A localized class may be divisible by \(2\); the next
lemma shows that such a division does not cost a factor at each switch.

\begin{lemma}[Exact determinant switch over a DVR]\label{rd:det-switch}
Let \(O\) be a DVR with fraction field \(F_O\).  Suppose that integral
perfect complexes \(C_{\rm str},C_F,C_S,C_{\rm rel}\), their
localization triangles, and their dualities are the base changes of
the diagram in Lemma~\ref{rd:local-model}.  Suppose the generic cohomology of these four global complexes is
concentrated in degrees one and two, that
\(H^1(C_F\otimes F_O)=F_OA\), and
\[
 f(A)=v_0\ne0,\qquad s(B)=Jv_0,\qquad f(B)=0
\]
for a class \(B\) in the switched complex.  The determinant tensors
are formed from each class and its functional under these same
dualities.  Then the switched generic cohomology has rank one,
is spanned by \(B\), and
\begin{equation}
 v_O(u_F)=v_O(u_S).
 \label{rd:valuation-switch}
\end{equation}
Neither localization is required to be primitive.
\end{lemma}

\begin{proof}
For this proof only, cohomology without a coefficient ring means
cohomology over \(F_O\).  Nonzero localization of \(A\) gives
\(H^1(C_{\rm str})=0\).  The relaxed localization image is
self-annihilating in \(F\oplus S\), so it has dimension two.
It contains the two independent vectors \((v_0,0)\) and
\((0,Jv_0)\) in the opposite planes.  Its intersections with either plane are therefore
exactly one-dimensional.  This proves the generic rank assertions.

Put \(V=H^1(C_{\rm rel})\), with ordered basis \(A,B\).
Duality identifies \(H^2(C_{\rm str})\) with \(V^*\); denote the
dual basis by \(A^*,B^*\).  The two integral triangles are
\begin{equation}
 C_{\rm str}\longrightarrow C_F\longrightarrow F[-1]\longrightarrow,
 \qquad
 C_{\rm str}\longrightarrow C_S\longrightarrow S[-1]\longrightarrow.
 \label{rd:triangles}
\end{equation}
Their fraction-field exact sequences are
\[
 0\longrightarrow F_OA\longrightarrow F
       \xrightarrow{\delta_F}V^*
       \longrightarrow F_OA^*\longrightarrow0
\]
and the same sequence with \(A,F,A^*\) replaced by \(B,S,B^*\).
The boundary is the dual of the relaxed localization.  In view of
\eqref{rd:cross-pairing} and \(s(B)=Jv_0\), it satisfies
\begin{equation}
 \delta_F(x)=\epsilon_F\langle x,v_0\rangle B^*,
 \qquad \epsilon_F\in O^\times.
 \label{rd:boundary}
\end{equation}
Choose \(x\in F\otimes F_O\) with \(\delta_F(x)=B^*\).
Then \(v_0\wedge x\) has unit volume in the integral determinant
lattice \(\det_O F\), because the alternating form on \(F\) is
unimodular and \(\langle x,v_0\rangle\) is a unit.  Although \(x\)
can be nonintegral, this assertion about its exterior product is
an equality of lattices.  Similarly choose \(y\in S\otimes F_O\)
mapping to \(A^*\).  The volume of \(Jv_0\wedge y\) is a unit in
\(\det_O S\).

The determinant isomorphisms of the two exact sequences send
\[
 A\otimes(A^*)^*
 \quad\hbox{and}\quad B\otimes(B^*)^*
\]
respectively, up to units and the irrelevant determinant sign, to
\begin{equation}
 (A^*\wedge B^*)^*\otimes(v_0\wedge x),
 \qquad
 (A^*\wedge B^*)^*\otimes(Jv_0\wedge y).
 \label{rd:common-tensor}
\end{equation}
These comparisons take place in
\(\mathcal D(C_{\rm str})\otimes\det F\) and
\(\mathcal D(C_{\rm str})\otimes\det S\).
The two plane determinants are identified by a unimodular map.
Hence the two coordinates have the same valuation.

This also explains exactly where the finite indices go.  For a
two-term elementary summand
\([O\xrightarrow{\pi^d}O]\) in degrees one and two, its contribution
to the coordinate in \(\mathcal D\) is \(\pi^{-d}\).
Such summands, including those arising from nonprimitive
localization, occur inside the integral isomorphisms
\eqref{rd:triangles}.  They are not discarded when writing
\eqref{rd:common-tensor}.  All strict torsion contributions are common
to the two sides.  The argument does not ask for an integral
complement to \(v_0\); for example \(v_0=2^de_1\) is paired with
the fractional vector \(2^{-d}e_2\), whose exterior product with
\(v_0\) has unit volume.
\end{proof}

\subsection{Rank reduction and the outer DVR}
\label{rd:outer-subsection}

The local switch preserves a determinant valuation.  We now use it to
reduce the dimension of any larger specialized Selmer space until a
nonzero minor controls that valuation.  The first argument handles each
characteristic-zero height-one prime.  A second, outer limit turns a
hypothetical sequence of unbounded losses at $(2)$ into the same rank
reduction problem.

\begin{lemma}[Selecting a rank-two evaluation]\label{rd:rank-two}
Let \(\mathcal G\) be a group, \(H\) a normal subgroup acting trivially
on an absolutely irreducible representation \(W\) over a field \(F_O\),
and suppose that restriction of the relevant cohomology classes to
\(H\) is injective.  For two independent classes, their joint
evaluations on \(H\) span \(W\oplus W\).  If the finite localization
is the rank-two projection supplied by \(1+\tau\) on the induced
representation, some actual evaluation has rank two.  One of the
two classes may be any prescribed nonzero class.
\end{lemma}

\begin{proof}
On \(H\) a crossed homomorphism is an additive homomorphism, and its
image is equivariant for \(\mathcal G\).  A proper submodule of
\(W\oplus W\) dominating both factors would impose a nonzero
scalar linear relation on the two evaluations: the direct sum of
two simple modules is semisimple, and
\(\operatorname{End}_{\mathcal G}(W)=F_O\).
That relation, followed by injectivity of restriction, would
contradict independence of the classes.  Thus the joint span is
the whole direct sum.  Projecting by \(1+\tau\) gives all pairs
in the two-dimensional finite plane.

It remains to pass from a span to a single element of the additive
evaluation image \(U\).  If the determinant quadratic polynomial
\(Q\) vanished on \(U\), then its polar form
\(Q(u+v)-Q(u)-Q(v)\) would vanish for \(u,v\in U\).
Consequently
\[
 Q\Bigl(\sum_i a_i u_i\Bigr)
 =\sum_i a_i^2 Q(u_i)
  +\sum_{i<j}a_i a_j
       \bigl(Q(u_i+u_j)-Q(u_i)-Q(u_j)\bigr)
 =0.
 \]
It would vanish on the entire field span, where it is not the zero
polynomial.  This contradiction proves the assertion.  This
polarization argument is valid in characteristic two as well.
\end{proof}

\begin{proposition}[Nonnegative valuations away from \((2)\)]
\label{rd:away-two}
Assume the simultaneous diagrams, derivative classes, and exact-kernel
evaluation maps of Lemmas~\ref{rl:evaluation}, \ref{rl:kernel},
\ref{rd:local-model}, and \ref{rd:derivative}.  For a primitive
constituent whose generic cohomology has rank at most one in degrees
one and two and zero elsewhere, its class-dual determinant
coordinate belongs to \(\Lambda[1/2]\).
This includes a simple-zero constituent and a base with
\(2^\infty\)-Selmer corank one.
\end{proposition}

\begin{proof}
The zero coordinate needs no argument.  Otherwise the generic
rank is one and the class spans it.  Localize at a height-one prime
\(\mathfrak p\ne(2)\), and denote the residue field by \(F_O\).
Its characteristic is zero.  Lemma~\ref{rl:local-acyclic}, the
degree-zero local conditions, and duality give zero fiber
cohomology in degrees zero and three.  Canceling unit blocks
therefore gives a free DVR model in degrees one and two.
A primitive vector in its generic kernel has nonzero image in its
\(\mathfrak p\)-fiber: the kernel of a map of free modules over a
DVR is saturated.  Include this vector in each localization test.
The two-sign reduction below has the same pattern as
\cite[Lemmas~1.5.3 and~1.6.4]{Howard2012}; the integral switch
and the evaluation hypotheses needed here at \(2\) have been
constructed above.

Suppose first that \(\Psi^2\ne1\) on \(F_O\).  Inertia at \(r\)
distinguishes the two summands of the induced representation.
The Tate representation over \(K\) is absolutely irreducible in
characteristic zero: this follows from the open Tate image, or
from its CM normalizer when \(K\) is not the CM field.
Thus the induced representation is absolutely irreducible.
If the test Selmer dimension is \(s>1\), Lemma~\ref{rd:rank-two}
and the finite Chebotarev realization in Lemma~\ref{rl:chebotarev-evaluation}
give a derivative prime with rank-two finite localization,
nonzero on the chosen primitive vector.

Poitou--Tate duality now determines the effect on the fiber.
The relaxed localization image is a self-annihilating plane.
It contains the whole finite plane, and hence equals that plane.
The strict kernel has dimension \(s-2\), and its replacement by
the singular condition adds no new class.  The fiber dimension
therefore drops by two.  On the generic field, the localization
is nonzero, so Lemmas~\ref{rd:derivative} and
\ref{rd:det-switch} preserve generic rank one, the nonzero
derivative class, and the determinant valuation.

If \(\Psi^2=1\), use conjugation, including its action
\((1+t)\mapsto(1+t)^{-1}\), to split the fiber into two signs.
All prior plane conditions are stable under this involution in
the split local model.  The two Galois summands are nonisomorphic
absolutely irreducible representations; the only CM quadratic
self-twist has been excluded.  Each sign has one finite and one
singular coordinate.  The local pairing in a sign is a
nondegenerate symmetric hyperbolic pairing, whose only isotropic
lines in characteristic zero are the finite and singular axes.
If both signs occur globally, choose one class from each, including
a nonzero component of the primitive vector, to obtain the same
rank-two reduction.

If only one sign occurs and \(s>1\), choose a rank-one finite
localization nonzero on that vector.  In this sign the relaxed
image is the finite line.  In the other sign its intersection
with the finite line is zero, and self-duality forces it to be
the singular line.  A switch retains the old strict kernel of
dimension \(s-1\ge1\), introduces the other sign, and preserves
the total dimension.  The next rank-two switch reduces it.
After finitely many such operations the fiber dimension is one.
It cannot fall below one, since a primitive vector from the
generic rank-one kernel survives after each switch.

At fiber dimension one, a minimal self-dual model over the DVR
has one free term in each of degrees one and two and zero
differential.  The class and paired functional are integral at
this height-one prime.  Their determinant coordinate is
therefore integral.  Reversing
\eqref{rd:valuation-switch} proves nonnegative valuation at
\(\mathfrak p\).  Applying this at every height-one prime other
than \((2)\), and using that the normal ring \(\Lambda\) is the
intersection of its height-one valuation rings inside its fraction
field, gives membership in \(\Lambda[1/2]\).

For the stated arithmetic instances, the rational central
specialization has Selmer dimension one and no invariants.
Upper semicontinuity bounds the generic dimension by one.
If the generic class is nonzero, it spans this space.  This
uses Selmer corank at an unknown base and does not replace that
corank by Mordell--Weil rank.
\end{proof}

\begin{lemma}[The terminal minor estimate]\label{rd:minor}
Let
\[
 C_i=[\,O_i^n\xrightarrow{D_i}O_i^n\,]
\]
be based complexes over DVRs with uniformizer \(2\), with generic
rank-one kernel.  Let \(z_i\in\ker D_i\) be integral, and let
\(\lambda_i\in 2^{-a}(O_i^n)^*\) annihilate \(\operatorname{im}D_i\),
where \(a\) is fixed.  If an \((n-1)\)-minor \(\Delta_i\) has
\(v(\Delta_i)\le b\), the determinant coordinate \(u_i\) of
\(z_i\otimes\lambda_i\) satisfies
\begin{equation}
 v(u_i)\ge-a-b.
 \label{rd:minor-bound}
\end{equation}
\end{lemma}

\begin{proof}
Permute rows and columns so that the minor is the upper-left block
\(A_i\).  Write \(D_i=\left(\begin{smallmatrix}A_i&b_i\\c_i&d_i
\end{smallmatrix}\right)\).  Since its rank is \(n-1\), its Schur
complement is zero.  Over the fraction field, use the triangular
basis changes with entries \(A_i^{-1}b_i,c_iA_i^{-1}\) to leave
\(A_i\) and a zero one-by-one block.  These triangular changes
have determinant one.  If \(\alpha_i\) is the last coordinate
of \(z_i\) and \(\beta_i\) the value of \(\lambda_i\) on the last
target basis vector, cancellation gives exactly
\begin{equation}
 u_i=\pm\frac{\alpha_i\beta_i}{\det A_i}.
 \label{rd:minor-formula}
\end{equation}
Here \(\alpha_i\in O_i\) and \(\beta_i\in2^{-a}O_i\).
This proves \eqref{rd:minor-bound}.  If \(n=1\), the empty
minor is \(1\) and the same formula applies.
\end{proof}

\begin{proposition}[Uniform residual valuation from outer rank reduction]
\label{rd:outer-bound}
Assume all of the following properties of a collection of primitive
ring instances:
\begin{enumerate}
\item after any fixed number of switches, minimal \(\mathcal O\)-models
have size bounded independently of \(h\) and of \(K\) with
\(\omega(k)\le A\);
\item the all-sequence evaluation and exact-kernel restriction of
Proposition~\ref{rl:outer-evaluation}, and the finite Chebotarev
realization of Lemma~\ref{rl:chebotarev-evaluation}, hold for these
models, retaining their degree-two cocycle identities;
\item the integral localization and duality diagrams and the
derivative classes of Lemma~\ref{rd:derivative} survive this same
outer limit; the class and its closed dual functional have a common
bounded clearing exponent.
\end{enumerate}
Then there is \(C(E,A)\) such that
\begin{equation}
 v_{(2)}(u)\ge-C(E,A)
 \label{rd:residual-bound}
\end{equation}
for every simple-zero ring instance.  The hypotheses refer to a
single simultaneous construction.  In particular, separate
abstract models realizing its individual cohomology groups do
not establish this proposition's arithmetic applicability.
\end{proposition}

\begin{proof}
The preceding paired-diagram construction supplies the common
clearing exponent in the arithmetic application.  The bound on
minimal sizes is Lemma~\ref{rl:residual-bound}; in particular,
active conductor primes do not contribute to that size after
residual localization.

Suppose that \(v_{(2)}(u_i)\to-\infty\).  All these coordinates
are nonzero.  First put the complexes into square two-term form.
At the residual \(\Omega\)-test, moving-prime inertia excludes
global invariants, and the local degree-zero conditions identify
Selmer invariants with those global invariants.  Thus \(H^0=0\).
The retained self-duality gives \(H^3=0\) and equal dimensions
in degrees one and two.  A minimal \(\mathcal O\)-model therefore
has only two free terms, of equal rank, in those degrees; all other
terms cancel in integral disks.  Fix a nonprincipal ultrafilter
\(\mathcal U\) and fix the bounded rank on an ultrafilter-large
set.  Write these models as
\[
 C_i=[\,O_i^n\xrightarrow{D_i}O_i^n\,].
\]

Transport the class, the uniformly cleared functional, all evaluation
matrices, and the finitely many retained comparison maps and boundary
witnesses through these integral disk contractions.  Their entries
remain integral and their identities remain exact.  The individual
active-place complexes, whose number may grow with the outer instance,
are not retained as an outer list of objects.  Their effect has already
been incorporated in the target cone and its determinant; only the
bounded minimal models and the fixed number of switch diagrams enter
this limit.

Apply Proposition~\ref{rl:outer-evaluation} to these models and
their retained integral diagrams.  It gives a DVR \(O_\infty\)
with uniformizer \(2\), characteristic-zero fraction field
\(F_\infty\), and injective all-sequence evaluation, including
restriction to the product of the exact inner kernels.
The induced representation uses one simultaneous complex
conjugation, so the sequence group over the quadratic fields has
index two in the rational sequence group.  Lemma~
\ref{rl:irreducibility} gives absolute irreducibility over
\(F_\infty\): the moving inertia character has infinite order,
and the finite Tate-matrix witness survives the outer quotient.

The outer complex may have more cohomology than any inner generic
complex.  Its new cycles can be represented by vectors \(w_i\)
for which \(v(D_iw_i)\to+\infty\), without exact inner lifts.
The retained degree-two evaluation matrices are integral, so the
evaluations of \(D_iw_i\) also tend to zero.  This is the part of
Proposition~\ref{rl:outer-evaluation} that makes these new classes
available for the following rank reduction.

Choose in each inner generic kernel a primitive vector \(v_i\).
Some coordinate is a unit on a set in \(\mathcal U\), so its
outer limit is nonzero.  Thus the outer kernel has dimension
at least one.  If that dimension is \(s>1\), apply
Lemma~\ref{rd:rank-two} to two outer classes, one the surviving
primitive vector.  Realize their rank-two evaluation by inner
derivative primes.  This realization only asks for finite
quotients at each stage.  For an approximate inner cocycle
one chooses the retained precision below the valuation of
its coboundary; no exact inner lift of the extra outer class
is being assumed.

More explicitly, retain the original inner index \(i\)
for outer instance \(j\), its ultrafilter \(\mathcal U_j\),
and the outer ultrafilter \(\mathcal U\).  Represent the
selected kernel element by \(u_j=(u_{j,i})_i\).
At stage \((j,i)\), apply finite Chebotarev to this actual
element and the old finite data to choose \(\ell_{j,i}\).
Adjoin its local complexes, the derivative cocycle and
the pairing homotopies to the same original stage
diagram with its old contractions.  Take the same
\(\mathcal U_j\) and then the same \(\mathcal U\).
The retained coefficient precision may increase more
slowly, but cofinally; the earlier maps and their
arbitrary evaluations are not reconstructed by an
independent choice of subsequence or model.

The localization of \(v_i\) is generically nonzero on a set
in \(\mathcal U\): its outer image is nonzero.  Therefore
the original inner spanning class also localizes nontrivially,
even if its multiple of \(v_i\) has large valuation.
Lemma~\ref{rd:det-switch} preserves each inner determinant
valuation exactly.  On the outer field, the relaxed image
is the finite plane and the singular switch lowers \(s\)
by two, by the same localization exact sequence as in
Proposition~\ref{rd:away-two}.  The new inner generic kernels
again supply primitive vectors, so the outer rank remains
at least one.

Only finitely many switches are necessary: the initial outer
rank is bounded and drops by two.  After each fixed number
of switches the required model-size bound remains available.
At the terminal outer rank one, the outer differential has
a nonzero maximal boundary minor.  Its inner determinants
\(\Delta_i\) have one finite valuation on a set in
\(\mathcal U\), by the definition of \(O_\infty\).
Lemma~\ref{rd:minor} then gives a fixed lower bound on
\(v(u_i)\), a contradiction.

This proof uses the minor only after rank reduction.  Bounded
matrix size alone would not suffice: the matrices \([2^i]\)
have bounded size and unbounded torsion length.  In that
example the entire matrix becomes zero in the outer quotient;
the rank-reduction and evaluation hypotheses would have to
supply additional arithmetic operations before a nonzero
terminal minor could be used.
\end{proof}

\subsection{Central specialization and the local indices}
\label{rd:specialization-subsection}

The preceding bounds concern the deformed determinant.  We next identify
its value at $t=0$ with the index of the original Heegner point and the
finite Shafarevich--Tate lengths.  The comparison retains the whole
integral complex; in particular, every active split conductor prime has
an explicit local contribution.

\begin{lemma}[Derived specialization at constant rational rank]
\label{rd:derived-specialization}
Let \(C\) be a perfect \(\Lambda\)-complex and put
\(C_0=C\otimes_\Lambda^{\mathbf L}\mathbb Z_2\), where \(t=0\).
Suppose its generic rational cohomology and its central rational
cohomology both have dimension one in degrees one and two and
are zero in other degrees.  Put \(A=\Lambda_{(t)}\), with residue field \(\mathbb Q_2\).
Let \(z\) be a degree-one cycle and \(\lambda\) a closed
degree-two functional in a finite free model of \(C\otimes_\Lambda A\),
and let \(z_0,\lambda_0\) be their images under \(A\to\mathbb Q_2\).
Then the generic determinant tensor \([z]\otimes[\lambda]\) is
regular at \((t)\) and specializes, under the determinant base-change
isomorphism, to \([z_0]\otimes[\lambda_0]\).  Its central valuation is
computed from the entire integral complex \(C_0\), including
its torsion cohomology.  It cannot be computed by substituting
the naive quotient of \(H^1(C)\) by \(t\).
\end{lemma}

\begin{proof}
Represent \(C\) by a bounded free complex.  Derived reduction
is its termwise reduction, so the determinant line itself
commutes with specialization.  For clarity, the cohomology
base-change sequence is
\begin{equation}
 0\longrightarrow H^i(C)/tH^i(C)
 \longrightarrow H^i(C_0)
 \longrightarrow H^{i+1}(C)[t]\longrightarrow0.
 \label{rd:tor-basechange}
\end{equation}
These last terms are the relevant \(\operatorname{Tor}_1\)
terms.  Diagonalize the complex over the DVR \(A\), or over its
completion \(\mathbb Q_2[[t]]\).  A nonunit pivot
would create an additional central cohomology class in its
two adjacent degrees.  The assumed constant dimensions
exclude such pivots.  Thus over \(\mathbb Q_2[[t]]\) there
are only invertible disks and the two free cohomology lines.
The rational class-dual tensor consequently specializes
to the stated central tensor.  Over \(\mathbb Z_2\), all
elementary divisors of \(C_0\), including the torsion supplied
by \eqref{rd:tor-basechange}, remain in its determinant
isomorphism.  This proves the assertion.
\end{proof}

\begin{lemma}[A ramified good odd place]\label{rd:active-local}
Let \(F_v/\mathbb Q_p\) be unramified, \(p\ne2\), and let
\(T'=T_2(E)\otimes\chi\), where \(E\) has good reduction
and \(\chi\) is ramified quadratic.  Then
\begin{equation}
 H^0(F_v,T')=0,\qquad
 H^1(F_v,T')=E^\chi(F_v)^{\wedge}_2,\qquad
 \operatorname{length}_{\mathbb Z_2}H^2(F_v,T')
 =\dim_{\mathbb F_2}E[2]^{\operatorname{Fr}_v}.
 \label{rd:local-length}
\end{equation}
The quotient of the full local complex by the integral Kummer
condition therefore contributes precisely this last length
in degree two.
\end{lemma}

\begin{proof}
The pro-\(2\) tame inertia generator acts by \(-1\).
Its complex has differential \(-2\), so its cohomology is
zero in degree zero and \(T'/2T'\) in degree one.
Taking residue-field cohomology gives
\[
 H^1(F_v,T')\simeq E[2]^{\operatorname{Fr}_v},\qquad
 H^2(F_v,T')\simeq E[2]/(\operatorname{Fr}_v-1)E[2].
 \]
The possible residue-cardinality factor on inertia cohomology
is one modulo \(2\).  Kernel and cokernel of an endomorphism
of a finite vector space have the same dimension, proving
the length assertion.

The integral local Kummer exact sequence identifies the
\(2\)-adic completion of the points with all of \(H^1(F_v,T')\).
Indeed the inverse Tate module of the local Weil--Chatelet
group is zero here: local duality identifies its
\(2\)-primary part with the dual of the finite completion
of the point group.  The point completion is finite because
the residue characteristic is odd.  Equivalently, inertia
\(-1\) implies that the local \(2\)-primary torsion is
exactly the Frobenius-fixed \(2\)-torsion.  The Kummer
complex includes the derived invariants of this finite
module on reduction, as in Lemma~\ref{rl:kummer}.
The quotient therefore has exactly the degree-two
cohomology displayed above.
\end{proof}

\begin{proposition}[The central determinant formula]
\label{rd:central-formula}
Suppose the product
\[
 L(E^{(h)},s)L(E^{(hk)},s)
\]
has a simple zero at \(s=1\).
Use the primitive Heegner sum \(P_h\) and its free-line index
\(j_h\) defined at the start of Section~\ref{rl:section}, and put
\[
 s(h)=\operatorname{length}_{\mathbb Z_2}
             \Sha(E^{(h)}/\mathbb Q)[2^\infty].
 \]
Use the primitive Selmer determinant coordinate \(u_h\)
with the fixed class multiplier of Lemma~\ref{rd:derivative}.
For the primitive Selmer diagrams constructed above, we have
\begin{equation}
 \begin{split}
 v(u_h(0))
  &=2j_h-
    \operatorname{length}_{\mathbb Z_2}
                  \Sha(E^{(h)}/K)[2^\infty]
    -2\sum_{p\mid h}t_p+O_E(1)\\
  &=2j_h-s(h)-s(hk)-4w(h)+O_{E,A}(1),
       \qquad \omega(k)\le A .
 \end{split}
 \label{rd:central-value}
\end{equation}
The first error includes the auxiliary prime and the fixed
parametrization multiplier.  The second has no dependence
on the number of primes dividing \(h\).
\end{proposition}

\begin{proof}
The two rational twists have analytic ranks zero and one.
The forward low-analytic-rank theorems give finiteness of
their whole Shafarevich--Tate groups; equivalently the
Heegner simple-zero theorem over \(K\) gives rank one
and finite Shafarevich--Tate group there.  Thus all finite
lengths in this proof refer to known finite groups.

The specialization comparison is made at each finite stage, before
the target data are retained.  Put \(c=|h|\).  At the split auxiliary
prime \(r_i\), the point distribution relation is
\[
 \operatorname{Tr}_{H_{cr_i}/H_c}(y_{cr_i})
       =(a_{r_i}-\sigma_{r_i}-\sigma_{\bar r_i})y_c,
\]
where the two Artin symbols have inverse ring-class actions.  At
\(t=0\), the scalar character is trivial, while the rational quadratic
character has value \(\chi_h(r_i)\) on both symbols.  The completed
weighted sum therefore gives exactly
\(L m_i P_h\), where \(m_i=a_{r_i}-2\chi_h(r_i)\).

This point identity is retained on bounded target complexes as follows.
Let \(C_h^{\rm Kum}\) be a fixed free model of the ordinary Kummer Selmer complex
over \(K\), and let \(b_h\) represent the Kummer class of \(P_h\).
At the retained coefficient precision, enlargement of the ramification
set and relaxation at the active and auxiliary primes give a target
comparison \(j_i:C_h^{\rm Kum}/2^i\to C_i/(t)\).  Each retained local
condition is an isomorphism on degree zero and an injection on degree
one, so the Selmer-cone exact sequence injects its degree-one
cohomology into global cohomology.  Kummer naturality and the point
relation therefore give, after contraction, a target degree-zero vector
\(v_i\) with
\[
 z_i(0)-L m_i j_i(b_h)=d v_i.
\]
Include this vector, the comparison matrix, and the associated local
localization triangles in the same target diagram.  Their ranks are
bounded for this fixed \(h,K\), independently of \(r_i\) and of
\([H_{cr_i}:K]\).  All statements are reduced to the same slower
cofinal precision if needed.  Lemma~\ref{rl:finite-limits} now retains
the displayed boundary identity.  The trace bound gives
\(v_2(m_i)\le c_E\); hence their limit \(m\in\mathbf Z_2\) on the
same ultrafilter is nonzero and has \(v_2(m)\le c_E\).
The active local quotients specialize
to Lemma~\ref{rd:active-local}, and the Kummer conditions at
\(2N\), including their differential lattices, are unchanged.
The central point is therefore nonzero, so generic
and central rational cohomology both have rank one.

To specialize the paired functional, choose its representative
over \(\Lambda_{(t)}\).  At an active ramified prime, inertia
minus one is \(-2\), a unit in this ring.  At \(r\), the
central Frobenius determinants are nonzero by the trace bound.
Their local complexes thus contract over \(\Lambda_{(t)}\).
Composing these contractions with retained global duality gives
the same generic functional as the contractions over
\(\mathcal O\), and specializes to rational central Kummer
duality.  Lemma~\ref{rd:derived-specialization} applies to this
representative; no regularity is claimed for an arbitrary
\(\mathcal O\)-cochain representing the functional.

The ordinary integral Kummer complex has compact Selmer
cohomology in degree one.  Its degree-two torsion has
length equal to that of the finite quotient of the
discrete Selmer group by its maximal divisible subgroup.
This follows by derived base change to
\(\mathbb Q_2/\mathbb Z_2\) and the exact local Kummer
conditions; the base-change terms must be retained.
Here that quotient is \(\Sha(E^{(h)}/K)[2^\infty]\).
The free parts of degrees one and two are paired
unimodularly by integral duality.  Torsion in degree
one and its dual contribution in degree three have
bounded length, by the fixed-scalar bound on global
torsion.

For completeness, an elementary divisor of size \(2^d\)
whose cokernel is in degree \(i\) contributes
\((-1)^{i+1}d\) to the coordinate in the inverse
determinant.  Consequently the degree-two
Shafarevich--Tate torsion subtracts its length.  On
the free line, a point of index \(2^{j_h}\) and its
functional each contribute \(j_h\), giving \(2j_h\).

At an active prime \(p\mid h\), both primes of \(K\)
above \(p\) are governed by
Lemma~\ref{rd:active-local}, each with length \(t_p\).
Relaxation therefore subtracts exactly \(2t_p\).
There is no bound hiding an error at every such prime.
At \(r\), the central scalar character is trivial;
the local degree-two determinant has the valuation
of \(r+1-a_r\chi_h(r)\) at each split place.
The trace condition
\[
 v_2(2-a_r),\ v_2(2+a_r)\le c_E,\qquad
 v_2(r-1)>c_E
 \]
makes both valuations at most \(c_E\).  It also bounds
the trace multiplier
\(a_r-2\chi_h(r)\) on the central point.
At primes over \(2N\) the integral Kummer condition
is retained.  The local twisting types form a fixed
finite set and \(K\) is split there; all comparison
indices, including those of Neron differential
lattices at \(2\), are therefore bounded in terms
of \(E\).  This proves the first equality, using
Lemma~\ref{rd:derived-specialization} to compare
the generic tensor with the central tensor.

We spell out the quadratic comparison needed for
the second equality.  Set
\[
 A_h=\operatorname{Res}_{K/\mathbb Q}(E^{(h)}_K),
 \qquad B_h=E^{(h)}\times E^{(hk)}.
 \]
There are isogenies in both directions whose
composites are multiplication by \(2\).
The isogeny-invariance identity for the BSD quotient,
with finite Shafarevich--Tate groups, is
\begin{equation}
 \frac{\#\Sha(A_h)\operatorname{Reg}(A_h)
             \operatorname{Per}(A_h)\prod_p c_p(A_h)}
      {\#A_h(\mathbb Q)_{\rm tors}
       \#A_h^\vee(\mathbb Q)_{\rm tors}}
 =
 \frac{\#\Sha(B_h)\operatorname{Reg}(B_h)
             \operatorname{Per}(B_h)\prod_p c_p(B_h)}
      {\#B_h(\mathbb Q)_{\rm tors}
       \#B_h^\vee(\mathbb Q)_{\rm tors}} .
 \label{rd:isogeny-identity}
\end{equation}
Regulators and periods are taken with compatible
dual and Neron normalizations.  This is the
Cassels--Tate isogeny identity.  The elliptic-curve local-measure
formula is due to Cassels \cite[Theorem~1.1]{Cassels1965}.
For the general abelian-variety identity used here, see the proof of
\cite[I, Theorem~7.3, especially (7.3.1), pp.~97--100]{MilneADT2006}
and \cite[Section~4.1]{DokchitserDokchitser2010}.  This identity requires
finite Shafarevich--Tate groups, as already established at these
noncentral twists, but does not assume the BSD conjecture.
The fixed-degree composites bound the rank-one
regulator indices, the torsion ratios, and the
differential and period indices by fixed powers
of \(2\).  The identity is essential: a statement
that the maps on Shafarevich--Tate groups have
kernel and cokernel killed by \(2\) would not
bound their lengths uniformly.

At every split prime of \(K/\mathbb Q\) the component
factors of \(A_h\) and \(B_h\) agree.  In particular
this cancels all primes dividing \(h\).
At an odd prime \(p\mid k\) outside \(2N\), the curve
over \(K\) has good reduction.  Its Weil restriction
has connected special fiber: the reduction of
Weil restriction over the ramified quadratic
extension is an extension of the good special
fiber by a connected additive group.  Thus its
component factor is \(1\).  Of the two rational
twists, one is good and the other is ramified
quadratic with good original curve.  The latter
has \(2\)-part of its component number equal to
\(2^{t_p}\).  This follows as well from the
ramified tame description in
Lemma~\ref{rd:active-local}: its connected
reduction is unipotent, and its local
\(2\)-primary torsion is the fixed
\(2\)-torsion.  At other good nonsplit places
the component factors are \(1\); at \(2N\)
there are only the fixed local comparisons.
Taking valuations in \eqref{rd:isogeny-identity}
therefore gives
\begin{equation}
 \operatorname{length}\Sha(E^{(h)}/K)[2^\infty]
 =s(h)+s(hk)+\sum_{p\mid k}t_p+O_E(1).
 \label{rd:sha-comparison}
\end{equation}
Since \(0\le t_p\le2\), the sum over \(k\) is at
most \(2A\).  Substituting this equality and
\(2\sum_{p\mid h}t_p=4w(h)\) proves the result.
\end{proof}

\begin{proof}[Proof of Corollary~\ref{rd:ring-lower}]
The arithmetic diagrams constructed with Lemma~\ref{rd:derivative},
their paired functionals, and the model bound of
Lemma~\ref{rl:residual-bound} give the hypotheses of
Propositions~\ref{rd:away-two} and \ref{rd:outer-bound}.
Those propositions give nonnegative height-one valuations
away from \((2)\) and a lower bound \(-C(E,A)\) at \((2)\).
Normality of \(\Lambda\) gives \(2^{C(E,A)}u_h\in\Lambda\).
Its constant coefficient has valuation at least \(-C(E,A)\).
Apply \eqref{rd:central-value}.
\end{proof}

\subsection{The universal correction and bounded clearing}
\label{rd:clearing-subsection}

For transfer to a missing vertex, the individual determinant coordinates
must belong to one quadratic group-ring family.  We first remove the
Euler factors at primes unused by a character.  We then clear the
remaining denominators with a factor whose central valuation and
negative Laurent support are bounded independently of the cube size.
The base curve and the companion field are fixed in this argument.

Fix \(K\) and a binary family \eqref{rd:binary-data}, based at
\(h_0=1\).  Put
\[
 G=(\mathbb Z/2\mathbb Z)^b,\qquad R=\Lambda[G].
\]
Let \(g_q\in G\) be the inertia bit:
\(g_q(x)=(-1)^{\lambda_q(x)}\).
Augmentation is denoted by subscript \(0\).

The universal class is constructed before making any character
evaluation.  Fix this finite family and \(G\).  At precision \(i\)
take the Heegner point of conductor \(r_i\prod_{q\in Q}q\).
Let \(\varepsilon:G_K\to G\) be the universal quadratic character.
In the finite Shapiro transfer of Lemma~\ref{rd:derivative}, use
the coefficient map
\[
 A_i[\operatorname{Gal}(H_{r_i\prod q}/K)]
       \longrightarrow A_i[G],\qquad
 [\sigma]\longmapsto\Psi_i(\sigma)[\varepsilon(\sigma)],
\]
with the same inverse-action convention on both sides.  Complete
this weighted coset sum to obtain a target Selmer-cone cycle over
\(A_i[G]\), with local lifts and boundary witnesses, and then
contract these target data.  The target-side rank bound applies
for this fixed family and gives one trace class over \(\Lambda[G]\).
Evaluation at \(x\) commutes with the finite coefficient map and
every target identity, so gives the full trace for that character.
If scalar models have already been chosen, compare the evaluated
model with them through the original target cochains and retain the
resulting comparison matrices and contraction homotopies.  There are
only finitely many characters for this fixed \(G\).  At \(t=0\), the
universal split-prime multiplier is
\[
 a_{r_i}-2[\varepsilon(\operatorname{Fr}_{\mathfrak r_i})].
\]
The conjugate Frobenius values agree because the character values
have exponent two.  Its evaluation at \(x\) is
\(a_{r_i}-2\chi_{h_x}(r_i)\), so the central boundary comparison
used in Proposition~\ref{rd:central-formula} is retained in this
same group-ring diagram before character evaluation.
It includes Euler factors at primes unused by \(x\), which the next
lemma removes.  No inverse limit over growing \(G\) or growing active
support is used in this construction.

The accompanying dual functional is constructed over
\(\mathcal O[G]\), not separately at its characters.  At a new
prime and at \(r\), the full local complex is acyclic modulo
the maximal ideal by Lemma~\ref{rl:local-acyclic}, hence has an
integral contraction over this local ring.  Use these contractions
to lift the class to the orthogonal Selmer complex, and then apply
the retained integral global duality.  This gives a closed integral
functional, up to one fixed old-support power of \(2\) when such
a support is included.  The present base is \(h_0=1\).

\begin{lemma}[Removing unused primes in the group ring]
\label{rd:unused}
Suppose the universal relaxed trace class and its duality
diagram have been constructed over \(R\).
There is \(\eta\in\mathcal O[G]\) such that
\(\eta u_{\rm rel}\) has each character coordinate
equal to the primitive determinant coordinate, up
to a \(\Lambda\)-unit.  The correction has no
denominator accumulating with \(\#Q\).
\end{lemma}

\begin{proof}
At an unused split prime \(q\), write \(z\) for the
full scalar Frobenius at one orientation.  Conjugation
gives \(z^{-1}\) at the other orientation.  The
point trace contributes
\[
 e_q=a_q-z-z^{-1},
 \]
whereas relaxation of the singular quotient contributes
\begin{equation}
 l_q=
 \det(1-q^{-1}z\operatorname{Fr}_q\mid T)
 \det(1-q^{-1}z^{-1}\operatorname{Fr}_q\mid T).
 \label{rd:unused-factors}
\end{equation}
Using reciprocal Frobenius conventions changes only
the corresponding unit conventions in the paired
determinant.  The localization triangle gives
\(u_{\rm rel}=e_q^2u_{\rm prim}/l_q\) at this
unused prime.  The square is required because the
point and its dual both occur.

Both factors are units in the local ring
\(\mathcal O[G]\).  To check this it suffices to
reduce modulo its maximal ideal, generated by \(2\)
and the augmentation ideal.  There the scalar
is \((1+t)^m\) with \(m\) odd, and the residual
Frobenius polynomial is nonzero at this nonconstant
scalar.  More precisely, before character
evaluation the following identity holds modulo \(2\):
\begin{equation}
 (1+a_qz+z^2)(1+a_qz^{-1}+z^{-2})
       =a_q^2+z^2+z^{-2}
       =(a_q+z+z^{-1})^2 .
 \label{rd:residual-euler}
\end{equation}
Thus \(l_q/e_q^2-1\in2\mathcal O[G]\), and
\begin{equation}
 \eta=
 \prod_{q\in Q}
 \left(1+
       \frac{l_q/e_q^2-1}{2}(1+g_q)\right)
 \in\mathcal O[G].
 \label{rd:eta}
\end{equation}
Its factor evaluates to \(l_q/e_q^2\) when \(q\)
is unused and to \(1\) when \(q\) is active.
This proves the claim.  The apparent inertia
projector denominator is canceled inside each
factor before taking the product.
\end{proof}

\begin{lemma}[Laurent expansion and simultaneous elimination]
\label{rd:schur}
There is a natural identification
\begin{equation}
 \mathcal O/2^M\mathcal O
       =(\mathbb Z/2^M\mathbb Z)((t)).
 \label{rd:laurent}
\end{equation}
Suppose \(A=\operatorname{diag}(A_1,\ldots,A_m)\)
has entries in \(\Lambda\), each residual determinant
is nonzero, and
\(\operatorname{ord}_t(\det\overline A_j)\le h\).
If \(B=A+2E\) with \(E\) over \(\Lambda\), then
every entry of \(B^{-1}\bmod2^M\) has no term below
\(t^{-Mh}\).  The same bound holds for a Schur
complement \(P-QB^{-1}S\), where \(P,Q,S\)
have power-series entries.  The bound is independent
of \(m\).
\end{lemma}

\begin{proof}
For a denominator \(s\in\Lambda\setminus(2)\),
write \(s=t^d v+2w\), with \(v\) a power-series
unit.  Its inverse modulo \(2^M\) is its finite
geometric expansion about \(t^dv\), which has
finite negative support.  Conversely all Laurent
series of finite negative support occur after
inverting \(t\).  This proves \eqref{rd:laurent}.
It also proves that an element of \(\mathcal O\)
has a unique series \(\sum a_nt^n\), \(a_n\in\mathbb Z_2\),
with \(v(a_n)\to+\infty\) as \(n\to-\infty\),
and that its valuation is \(\min_n v(a_n)\).

Write \(\det A_j=t^{d_j}v_j+2w_j\), with
\(d_j\le h\), and put
\[
 C_j=\frac{\operatorname{adj}(A_j)}{t^{d_j}v_j},
 \qquad C=\operatorname{diag}(C_j).
 \]
Each entry of \(C\) has pole at most \(h\), and
\(BC=1+2F\), where entries of \(F\) have the same
pole bound.  Therefore
\begin{equation}
 B^{-1}\equiv C\sum_{j=0}^{M-1}(-2F)^j
                                  \pmod{2^M}.
 \label{rd:inverse-expansion}
\end{equation}
The \(j\)-th summand has pole at most \((j+1)h\).
Adding more entries in a matrix product cannot
worsen this pole bound.  Multiplication by
power-series matrices proves the Schur assertion.

For a block over \(R\) whose augmentation is \(A\),
the block is invertible over \(\mathcal O[G]\):
its residual determinant in the local ring is a
unit.  Each sign evaluation is \(A+2E_x\) with
power-series entries, so the same estimate holds
at all characters.  If \(S_0=0\), its Schur
complement has augmentation exactly \(P_0\).
\end{proof}

\begin{lemma}[Clearing a bounded complex]\label{rd:clearing}
Let \(D\) be a based complex of finite free
\(\mathcal O[G]\)-modules of total graded rank
at most \(r_0\), with \(D_0\) defined over
\(\Lambda\).  Suppose \(D_0\otimes_\Lambda^{\mathbf L}\mathbb Z_2\)
has rational cohomology of dimension one in degrees
one and two, zero elsewhere, and total integral
torsion length at most \(h_1\).
Let \(z\) be an integral cycle and let
\(\lambda\in2^{-a_0}(D^2)^*\) be a closed functional.
Assume every character constituent is either
generically acyclic or has the same generic rank-one pattern;
in the generically acyclic case its determinant coordinate
is defined to be zero.  Multiplication of these
coordinates by any \(\eta\in\mathcal O[G]\)
is permitted.  Then there are \(f,d\in\mathcal O[G]\)
such that
\begin{equation}
 f_xu_x=d_x,\qquad
 f_0\in\Lambda,\qquad
 v(f_0(0))\le C(r_0,h_1,a_0)<\infty.
 \label{rd:clearing-equations}
\end{equation}
If all entries of \(D_x\) have bounded negative
Laurent support at each precision, uniformly
in \(b\), the same is true of \(f_x\).
No such bound on \(d_x,z,\lambda,\eta\) is required.
\end{lemma}

\begin{proof}
Put the central integral complex into elementary-divisor
bases over \(\mathbb Z_2\).  Its positive pivot
valuations sum to at most \(h_1\).  Lift these
constant integral basis changes.  Formally cancel
the same pivot positions upstairs, recording the
pivot numerators and all denominators in the
resulting rational basis changes and determinant
comparison.  There are at most \(r_0\) such
operations.  A product \(P\) of sufficiently
large, but bounded, powers of their numerators
clears every denominator.  Its degree in the
differential entries is bounded in terms of
\(r_0\); its augmentation belongs to \(\Lambda\)
and has central valuation bounded in terms of
\(r_0,h_1\).

On constituents where \(P_x\ne0\), the remaining
complex is one-by-one in degrees one and two.
If its differential is nonzero, the cycle
identity forces the projected cycle to be zero,
so the same rational determinant expression
gives the stipulated zero coordinate.  Otherwise
the expression is the product of the projected
cycle and functional, with the canceled-pivot
factor.  Clearing denominators gives
\[
 u_x=2^{-a_0}\eta_x Q_x/P_x,\qquad
 Q\in\mathcal O[G].
 \]
Take \(f=2^{a_0}P^2\) and \(d=\eta P Q\).
On a constituent where a pivot first vanishes,
\(P_x=0\), so both sides of \(f_xu_x=d_x\)
are zero.  Thus the identity holds on every
character.  Since \(f\) is a polynomial of
bounded degree in the bounded differential
entries, its central valuation and pole
properties follow.  All unused-prime
corrections and eliminated-block units occur
in \(d\); none is inserted into \(f\).
\end{proof}

\begin{proposition}[The bounded clearing factor for the family]
\label{rd:family-clearing}
Assume the compatible universal Selmer and Heegner
diagrams and their primitive specializations.
Suppose every character constituent of the universal Selmer complex
is generically acyclic or has generic
cohomology of dimension one in degrees one and two and zero
elsewhere.  If the fixed base has \(2^\infty\)-Selmer corank
one over \(K\), the corrected coordinates of
Lemma~\ref{rd:unused} admit the factorization
\eqref{rd:clearing-equations}, with central
valuation and fixed-precision pole bounds
independent of \(b\) and \(\#Q\).
\end{proposition}

\begin{proof}
Use the integral universal class and paired functional just
constructed.  To apply Lemma~\ref{rd:clearing}, it remains to
reduce the complex to bounded size without losing uniform control
of its negative Laurent powers.

At augmentation, the localization triangles
present the complex as the old complex extended
by a direct sum of singular local blocks in
degrees one and two, one block for each new
prime.  Its differential is upper triangular,
with the old complex a subcomplex.  This graded
model lifts over \(R\): split disks over the
augmentation local ring, identify the minimal
parts, and lift the invertible basis changes.
This is an application of the actual finite
diagram construction, rather than just an
equality of Euler characteristics.

Each new block has determinant \(l_{q,0}\).
Its residual \(t\)-order is uniformly bounded.
In the reducible residual range the residual
Frobenius polynomial is \((X+1)^2\), and an
odd scalar exponent gives order two at each
split place; thus order four suffices for the
two-place block.  Eliminate the entire new
block over \(\mathcal O[G]\).  By
Lemma~\ref{rd:schur}, the resulting complex
\(D\) has the bounded old size, has augmentation
equal to the old complex, and has a uniform
pole bound at every fixed precision.

The central cohomology of the old complex is
the fixed base Kummer cohomology, modified only
by the bounded auxiliary-prime terms.  Its
rational dimension is one in degrees one and
two.  The quotient of the discrete Selmer
group by its maximal divisible subgroup is
finite because its Pontryagin dual is a
finitely generated \(\mathbb Z_2\)-module.
For this fixed \(K,E\), its length is a finite
constant, even when the Shafarevich--Tate
group itself is not known to be finite.
Together with global torsion and the trace
bound at \(r\), this bounds all central
torsion lengths of \(D_0\).

Apply Lemma~\ref{rd:clearing}.  The determinant
of the eliminated block is a unit over
\(\mathcal O[G]\), so its inverse is an
allowed integral multiplier.  Include it
and \(\eta\) from \eqref{rd:eta} in the
numerator multiplier of that lemma.
This proves the assertion.

The dependence here is on fixed \(K,E\);
it is different from the uniform
\(C(E,A)\) of Proposition~\ref{rd:outer-bound}.
No uniform bound on the finite Selmer quotient
as \(K\) varies is asserted or needed for
the missing-vertex argument.
\end{proof}

\begin{proposition}[Assembly of the arithmetic diagrams]
\label{rd:assembly}
For the ring data stated at the start of this section, the actual
finite Heegner and Selmer constructions satisfy the simultaneous
diagram hypotheses of Propositions~\ref{rd:away-two},
\ref{rd:outer-bound}, and \ref{rd:central-formula}.
The bounds needed for the lower-bound argument depend only on
\(E\), the fixed local data, and a bound \(A\) on \(\omega(k)\).
For the binary data \eqref{rd:binary-data} with fixed \(E,K\)
and base Selmer corank one, assume in addition that every
character constituent of the universal Selmer complex is generically acyclic
or has generic cohomology of
dimension one in degrees one and two and zero elsewhere.
Then the constructions satisfy all the hypotheses of
Proposition~\ref{rd:family-clearing}, with
bounds independent of the cube dimension and of its
new-prime support.
\end{proposition}

\begin{proof}
The paired arithmetic diagram constructed after
Lemma~\ref{rd:derivative} contains the actual weighted target classes,
local Kummer lifts, target comparison and boundary witnesses, and
simultaneous duality triangles.  The target-side construction
gives its inner rank bound for each fixed finite conductor support;
no ring-class source complex is passed to the limit.
Theorem~\ref{rl:limit-evaluation-package} retains
their evaluations and exact-kernel restrictions.  The same-stage
extension in Proposition~\ref{rd:outer-bound} shows explicitly
that these data remain compatible after each further switch.

Lemma~\ref{rl:residual-bound} bounds the minimal models in terms
of \(E,A\) and the number of retained derivative primes, not the
active conductor support.  Outer rank reduction uses only a bounded
number of switches, hence only finitely many of these bounds.
The paired functional has the common clearing exponent established
above.  These are the three hypotheses of
Proposition~\ref{rd:outer-bound}; the characteristic-zero
evaluations and local acyclicity give those of
Proposition~\ref{rd:away-two}.  The actual central trace and
\(\Lambda_{(t)}\)-representative, together with the local
quotients of Lemma~\ref{rd:active-local}, were verified in the
proof of Proposition~\ref{rd:central-formula}.

For a binary family, Section~\ref{rd:clearing-subsection}
constructs the universal trace and its dual functional over the
group ring before evaluating characters.  Lemma~\ref{rd:unused}
then removes the unused-prime factors integrally.  The bounded
complex in Proposition~\ref{rd:family-clearing} comes from the
actual localization triangles and simultaneous Schur elimination,
with the order-four residual local bound.  Its central torsion is
controlled by the finite quotient of the fixed base Selmer group
and the bounded auxiliary-prime terms.  Together with the stated
constituent-rank assumption, these give all the hypotheses of
Lemma~\ref{rd:clearing}.

For varying \(K\), only its bounded ramification count enters
the residual-model and local-comparison constants.  For the final
binary argument \(K\) is fixed, so its finite base Selmer
quotient may enter the clearing constants.  At \(2N\) the
original Kummer conditions have been retained throughout; the
construction imposes no ordinary or semistable hypothesis.
\end{proof}

\subsection{Coefficient comparison and the missing vertex}
\label{rd:transfer-subsection}

The bounded clearing factor reduces the last step to finitely many
coefficient congruences.  A parity argument on a sufficiently large
binary cube makes a nonzero address agree with the base to the required
precision.  If the base derivative were zero, that congruence would
contradict the upper bound already supplied at every nonzero address.

The common-zero count in the following proof is the characteristic-two
Chevalley--Warning principle
\cite[Section~1, equation~(1)]{HeathBrownChevalleyWarning}; its short parity
proof is included here.

\begin{lemma}[A nonzero address with prescribed coefficient bits]
\label{rd:bits}
Let \(A_1(x),\ldots,A_q(x)\) be specified Laurent
coefficients of sign evaluations of elements of
\(\mathcal O[G]\).  For \(M\ge1\), if
\begin{equation}
 b>q(2^M-1),
 \label{rd:bit-threshold}
\end{equation}
there is \(x\ne0\) with
\(A_i(x)\equiv A_i(0)\pmod{2^M}\) for every \(i\).
\end{lemma}

\begin{proof}
Write the element supplying the \(i\)-th test as
\(a_i=\sum_{g\in G}a_{i,g}(t)g\), and let \(n_i\) be the Laurent
degree being tested.  Lemma~\ref{rd:schur} gives integral Laurent
coefficients, so
\[
 z_i=\sum_{g\in G}[t^{n_i}]a_{i,g}(t)\,g\in\mathbb Z_2[G].
\]
The group \(G\) is finite, hence coefficient extraction commutes
with character evaluation: \((z_i)_x=A_i(x)\).  Apply
Lemma~\ref{cy:binary} to these \(q\) integral group-ring elements.
Its bound is exactly \eqref{rd:bit-threshold}; no bound on the
Laurent degrees or on the coefficients enters this argument.
\end{proof}

\begin{lemma}[Transfer through the bounded factor]
\label{rd:binary-transfer}
Suppose \(f,d\in\mathcal O[G]\),
\(u_x\in\Lambda[1/2]\), and \(f_xu_x=d_x\)
for every character.  Suppose
\(f_0\in\Lambda\) and
\(c=v(f_0(0))<\infty\).
If every \(f_x\bmod2^M\) has no powers below
\(-L(M)\), then for
\[
 M>c,\qquad
 b>(L(M)+2)(2^M-1)
 \]
there is \(x\ne0\) satisfying
\begin{equation}
 v\bigl(u_x(0)-u_0(0)\bigr)\ge M-2c.
 \label{rd:binary-congruence}
\end{equation}
\end{lemma}

\begin{proof}
Apply Lemma~\ref{rd:bits} to the coefficients
of \(f_x\) in degrees \(-L(M),\ldots,0\)
and the constant coefficient of \(d_x\).
At its nonzero solution, all negative
coefficients of \(f_x\) are divisible by
\(2^M\), and its constant coefficient is
congruent to \(f_0(0)\).  Hence the latter
has valuation \(c\), and
\(v_{\mathcal O}(f_x)\le c\).
Since \(d_x\) is integral over \(\mathcal O\),
\(v_{\mathcal O}(u_x)\ge-c\), and the same
holds for \(u_0\).  The Laurent expansion
in Lemma~\ref{rd:schur} then shows that all
their coefficients lie in \(2^{-c}\mathbb Z_2\).
They have no negative coefficients because
they belong to \(\Lambda[1/2]\).

The constant-term convolution converges
\(2\)-adically, and gives
\[
 [t^0]d_x\equiv [t^0]f_x\,u_x(0)
                          \pmod{2^{M-c}},
 \qquad
 [t^0]d_0=f_0(0)u_0(0).
 \]
Use the imposed congruences and divide by
\(f_0(0)\).  This proves
\eqref{rd:binary-congruence}.  The possibly
large negative support of \(d_x\) was never
tested; only its constant coefficient was used.
\end{proof}

\begin{proof}[Proof of Theorem~\ref{rd:missing-vertex}]
For every character the primitive generic
cohomology has rank at most one.  At nonzero
vertices this follows from the simple-zero
theorem; at zero it follows from
\eqref{rd:base-corank} and rational
specialization.  There are no generic invariants, and duality
pairs degrees one and two; thus each primitive constituent is acyclic
or has exactly the rank-one pattern.  The unused-prime local blocks
are acyclic over the fraction field, by the nonzero factors
\eqref{rd:unused-factors}.  The universal Selmer complex has the
same generic pattern, as required by
Proposition~\ref{rd:family-clearing}.  With this hypothesis
verified, Proposition~\ref{rd:assembly} supplies the compatible
universal arithmetic construction.  By Proposition~\ref{rd:away-two}
and Lemma~\ref{rd:unused}, the corrected
coordinates belong to \(\Lambda[1/2]\).
Proposition~\ref{rd:central-formula} gives
\[
 v(u_x(0))\le B_1\quad(x\ne0)
 \]
for a constant \(B_1\) depending only on
\(B,E,K\) and the fixed local data.

If the base derivative vanished, the
Gross--Zagier formula would make the
central Heegner sum torsion: its canonical
height is a fixed nonzero multiple of
that derivative.  Here only the conductor-one formula at
the fixed base \(h=1\) is used, with the corrected primitive
sum and the split Heegner data of
\cite[Theorem~1.1]{CaiShuTian2014}.
At the base this implies \(u_0(0)=0\).
Indeed, if the generic rank is zero the
coordinate was defined to be zero.  If the
generic rank is one, its central rank is
also one by \eqref{rd:base-corank}, so
Lemma~\ref{rd:derived-specialization}
identifies its specialized class with
the zero rational Kummer class of that
torsion point.

Proposition~\ref{rd:family-clearing} supplies
\(f,d\) with a common bound \(c_*\) on
\(v(f_0(0))\) and common pole bounds
\(L(M)\).  Choose
\[
 M>\max\{c_*,B_1+2c_*\}
 \]
and then take \(b\) above the threshold
of Lemma~\ref{rd:binary-transfer}.
That lemma gives a nonzero vertex with
\(v(u_x(0))\ge M-2c_*>B_1\), a
contradiction.  Thus the base derivative
is nonzero.  The classical split Heegner
sign is odd, so the base product has
a simple zero.
\end{proof}

\begin{remark}[The simultaneous arithmetic interface]
\label{rd:interface-status}
The determinant argument can be reused only with
the simultaneous arithmetic construction verified
in Proposition~\ref{rd:assembly} and
Theorem~\ref{rl:limit-evaluation-package}.
Neither the published duality theorem nor the
classical Heegner formulas alone supplies that
interface.
\end{remark}
